\documentclass[11pt,reqno]{amsart}

\usepackage[T1]{fontenc}
\usepackage{geometry}
\geometry{margin=1in}

\usepackage{amsmath,amssymb,amsthm,mathtools}
\usepackage{microtype}
\usepackage{booktabs}
\usepackage[shortlabels]{enumitem}
\usepackage[unicode,colorlinks=true,linkcolor=blue,urlcolor=magenta,citecolor=blue]{hyperref}
\usepackage[capitalize,nameinlink]{cleveref}

% Paired delimiters and differential
\DeclarePairedDelimiter{\abs}{\lvert}{\rvert}
\DeclarePairedDelimiter{\norm}{\lVert}{\rVert}
\DeclarePairedDelimiterX{\inner}[2]{\langle}{\rangle}{#1, #2}
\newcommand{\dd}{\mathop{}\!\mathrm{d}}

% Common operators and spaces
\DeclareMathOperator{\supp}{supp}
\DeclareMathOperator{\dist}{dist}
\DeclareMathOperator{\pv}{p.v.}
\newcommand{\R}{\mathbb{R}}
\newcommand{\N}{\mathbb{N}}
\newcommand{\flap}{(-\Delta)^{s}}

% Theorem environments (AMS numbering within sections)
\theoremstyle{plain}
\newtheorem{theorem}{Theorem}[section]
\newtheorem{lemma}[theorem]{Lemma}
\newtheorem{proposition}[theorem]{Proposition}
\newtheorem{corollary}[theorem]{Corollary}

\theoremstyle{definition}
\newtheorem{definition}[theorem]{Definition}
\newtheorem{assumption}[theorem]{Assumption}
\newtheorem{example}[theorem]{Example}

\theoremstyle{remark}
\newtheorem{remark}[theorem]{Remark}

\numberwithin{equation}{section}

\title[Short Running Title]{Full Paper Title}
\author{Author Name}
\address{Department, Institution}
\email{author@example.edu}
\date{\today}

\begin{document}

\begin{abstract}
Concise summary of the problem, main theorem, and key analytic techniques.
\end{abstract}

\maketitle

\section{Introduction}\label{sec:intro}

State the mathematical setting, motivate the main objects, and summarize the primary theorem.
For $s \in (0, 1)$ and $d \ge 1$, let $u \in \mathcal{S}(\R^d)$ be a Schwartz function.

\begin{definition}[Fractional Laplacian]\label{def:flap}
For $s \in (0, 1)$ and $u \in \mathcal{S}(\R^d)$, the fractional Laplacian is defined by
\begin{equation}\label{eq:flap_pv}
  \flap u(x) \coloneqq c_{d,s} \pv \int_{\R^d} \frac{u(x) - u(y)}{\abs{x - y}^{d + 2s}} \dd y,
\end{equation}
where $c_{d,s} > 0$ is the standard normalization constant.
\end{definition}

\begin{theorem}[Energy Identity]\label{thm:energy}
Let $s \in (0, 1)$ and suppose that $u \in C([0, T]; L^2(\R^d)) \cap L^2(0, T; \dot{H}^s(\R^d))$
solves $\partial_t u + \flap u = 0$ on $\R^d \times (0, T)$. Then, for all $t \in [0, T]$,
\begin{equation}\label{eq:energy_id}
  \frac{1}{2}\norm{u(t)}_{L^2(\R^d)}^2
  + \int_0^t \norm{(-\Delta)^{s/2} u(\tau)}_{L^2(\R^d)}^2 \dd\tau
  = \frac{1}{2}\norm{u(0)}_{L^2(\R^d)}^2.
\end{equation}
\end{theorem}

\begin{proof}
Testing the equation against $u$ and integrating in time over $[0, t]$ yields
\begin{align*}
  \frac{1}{2}\norm{u(t)}_{L^2(\R^d)}^2 - \frac{1}{2}\norm{u(0)}_{L^2(\R^d)}^2
  &= -\int_0^t \inner{\flap u(\tau)}{u(\tau)}_{L^2(\R^d)} \dd\tau \\
  &= -\int_0^t \norm{(-\Delta)^{s/2} u(\tau)}_{L^2(\R^d)}^2 \dd\tau. \qedhere
\end{align*}
\end{proof}

\end{document}
